\usepackage{tikz}
\usepackage{pdfpages}
\usepackage{xcolor}
\usepackage{float}
\usepackage{caption}
\usetikzlibrary{fit, calc, positioning, matrix, arrows.meta, shadows}

% Вариант 1: Создание команды \highlight
\newcommand{\highlight}[1]{\colorbox{yellow!50}{#1}}

\renewcommand{\makethesistitle}[7]
{
	\documentmeta{РПЗ к ВКР}{#4}{}{#3}
	
	\begin{titlepage}
		\centering
		
		\titlepageheader{#1}{#2}
		\vspace{15.8mm}
		
		\titlepagenotetitle{К ВЫПУСКНОЙ КВАЛИФИКАЦИОННОЙ РАБОТЕ}{#3}
		\vfill
		
		\titlepageauthors{#4}{Руководитель ВКР}{#5}{#6}{#7}
		\vspace{14mm}
		
		\text{{\the\year} г.}
	\end{titlepage}
	
	\setcounter{page}{2}
}

% Создание шапки титульной страницы
\renewcommand{\titlepageheader}[2]
{
	\begin{wrapfigure}[7]{l}{0.14\linewidth}
		\vspace{3.4mm}
		\hspace{-8mm}
		\includegraphics[width=0.89\linewidth]{bmstu-logo}
	\end{wrapfigure}
	
	{
		\singlespacing \small
		Министерство науки и высшего образования Российской Федерации \\
		Федеральное государственное автономное образовательное учреждение \\
		высшего образования \\
		<<Московский государственный технический университет \\
		имени Н.Э.~Баумана \\
		(национальный исследовательский университет)>> \\
		(МГТУ им. Н.Э.~Баумана) \\
	}
	
	\vspace{-4.2mm}
	\vhrulefill{0.9mm} \\
	\vspace{-7mm}
	\vhrulefill{0.2mm} \\
	\vspace{2.8mm}
	
	{
		\small
		ФАКУЛЬТЕТ ИУ \longunderline{<<#1>>} \\
		\vspace{3.3mm}
		КАФЕДРА ИУ7 \longunderline{<<#2>>} \\
	}
}

% Установка заголовков РПЗ
\renewcommand{\titlepagenotetitle}[2]
{
	{
		\LARGE \bfseries
		РАСЧЕТНО-ПОЯСНИТЕЛЬНАЯ ЗАПИСКА \\
	}
	\vspace{5mm}
	{
		\Large \itshape
		#1 \\
		\vspace{5mm}
		НА ТЕМУ: \\
		\vspace{5mm}
		<<#2>> \\
	}
}

% Создание страницы определений
\renewenvironment{definitions}
{
	\chapter*{ОПРЕДЕЛЕНИЯ}
	\addcontentsline{toc}{chapter}{ОПРЕДЕЛЕНИЯ}
	
	В настоящей расчетно-пояснительной записке применяют следующие термины с соответствующими определениями.
	
	\begin{description}[leftmargin=0pt]
	}
	{
	\end{description}
}

% Создание страницы обозначений и сокращений
\renewenvironment{abbreviations}
{
	\chapter*{ОБОЗНАЧЕНИЯ И СОКРАЩЕНИЯ}
	\addcontentsline{toc}{chapter}{ОБОЗНАЧЕНИЯ И СОКРАЩЕНИЯ}
	
	В настоящей расчетно-пояснительной записке применяют следующие сокращения и обозначения.
	
	\begin{description}[leftmargin=0pt]
	}
	{
	\end{description}
}

\RequirePackage{enumitem}
\newcommand{\abbreviation}[2]
{
	\item \noindent #1 --- #2
}

\newcommand{\ssr}[1]
{
	\chapter{#1}
	\addcontentsline{toc}{chapter}{#1}  
}

\newcommand{\ssrno}[1]
{
	\chapter*{#1}
	\addcontentsline{toc}{chapter}{#1}  
}


% Переопределение формата даты обращения
\DeclareFieldFormat{urldate}
{%
	(дата~обращения:%
	\addspace
	\iffieldundef{urlday}{}{\ifnumless{\thefield{urlday}}{10}{0\thefield{urlday}}{\thefield{urlday}}%
		\adddot
		\iffieldundef{urlmonth}{}{\ifnumless{\thefield{urlmonth}}{10}{0\thefield{urlmonth}}{\thefield{urlmonth}}%
			\adddot
			\thefield{urlyear}}})
}

% Переопределение списков
\setlist[itemize]{label=---}

% Команды для различных типов слоев
\newcommand{\convlayer}[4][]{
	% Сверточный слой
	\foreach \i in {0,...,#4} {
		\node[layer, convcolor, right=2cm of #2, yshift=\i*(-6pt)+(#4*3pt), xshift=\i*(-6pt)+(#4*3pt), #1] (#3\i) {};
	}
}

\newcommand{\poollayer}[4][]{
	% Подвыборочный слой
	\foreach \i in {0,...,#4} {
		\node[smalllayer, poolcolor, right=2.5cm of #2, yshift=\i*(-6pt)+(#4*3pt), xshift=\i*(-6pt)+(#4*3pt), #1] (#3\i) {};
	}
}

\newcommand{\fclayer}[4][]{
	% Полносвязный слой
	\foreach \i in {0,...,#4} {
		\node[fclayer, fccolor, right=2.75cm of #2, yshift=\i*(-6pt)+(#4*3pt), xshift=\i*(-6pt)+(#4*3pt), #1] (#3\i) {};
	}
}

% Команда для капсульных слоев
\newcommand{\capsulelayer}[4][]{
	% Капсульный слой
	\foreach \i in {0,...,#4} {
		\node[layer, capsulecolor, right=2.5cm of #2, yshift=\i*(-6pt)+(#4*3pt), xshift=\i*(-6pt)+(#4*3pt), #1] (#3\i) {};
	}
}

% Команда для капсульных слоев
\newcommand{\attentionlayer}[4][]{
	% Капсульный слой
	\foreach \i in {0,...,#4} {
		\node[layer, attentioncolor, right=2.5cm of #2, yshift=\i*(-6pt)+(#4*3pt), xshift=\i*(-6pt)+(#4*3pt), #1] (#3\i) {};
	}
}

\newcolumntype{Y}{>{\centering\arraybackslash}X}

\makeatletter
% === Новое определение \@dottedtocline (без выключки по ширине) ===
\def\@dottedtocline#1#2#3#4#5{%
	\ifnum #1>\c@tocdepth \else
	\vskip \z@ \@plus.2\p@
	{\leftskip #2\relax
		\rightskip 0pt plus 1fil         % <--- левый край: правый отступ с бесконечной растяжимостью
		\parfillskip 0pt                 % <--- запрет растягивания последней строки
		\parindent #2\relax\@afterindenttrue
		\interlinepenalty\@M
		\leavevmode
		\@tempdima #3\relax
		\advance\leftskip \@tempdima \null\nobreak\hskip -\leftskip
		{#4}\nobreak
		\leaders\hbox{$\m@th \mkern \@dotsep mu\hbox{.}\mkern \@dotsep mu$}\hfill
		\nobreak
		\hb@xt@\@pnumwidth{\hfil\normalfont \normalcolor #5%
			\kern-\p@\kern\p@}%
		\par}%
	\fi}

% === Главы (с полужирным шрифтом) ===
\renewcommand*{\l@chapter}[2]{%
	\ifnum \c@tocdepth>\m@ne
	\addpenalty{-\@highpenalty}
	\vskip 1em \@plus.2em
	\@dottedtocline{0}{0pt}{1.5em}{\bfseries #1}{\bfseries #2}%
	\fi}

% === Разделы ===
\renewcommand*{\l@section}[2]{%
	\ifnum \c@tocdepth>\z@
	\addpenalty\@secpenalty
	\vskip 0 \@plus\p@
	\@dottedtocline{1}{1.5em}{2.3em}{#1}{#2}%
	\fi}

% === Подразделы ===
\renewcommand*{\l@subsection}[2]{%
	\ifnum \c@tocdepth>\@ne
	\addpenalty\@secpenalty
	\vskip 0 \@plus\p@
	\@dottedtocline{2}{3.8em}{3.2em}{#1}{#2}%
	\fi}
\makeatother

% === Содержание с запретом переносов ===
\renewcommand{\maketableofcontents}
{
	\renewcommand\contentsname{СОДЕРЖАНИЕ}
	\bgroup
	\hyphenpenalty=10000
	\exhyphenpenalty=10000
	\tableofcontents
	\egroup
}

%% bmstu_defabbr.sty
%% Copyright 2023 Mikael Novikov
%
% This work may be distributed and/or modified under the
% conditions of the LaTeX Project Public License, either version 1.3
% of this license or (at your option) any later version.
% The latest version of this license is in
%   http://www.latex-project.org/lppl.txt
% and version 1.3 or later is part of all distributions of LaTeX
% version 2005/12/01 or later.
%
% This work has the LPPL maintenance status `maintained'.
% 
% The Current Maintainer of this work is Mikael Novikov.
%
% This work consists of all files listed in manifest.txt.

% Создание элемента
\RequirePackage{enumitem}
\renewcommand{\definition}[2]
{
	\item[---] \noindent #1 -- #2
}

% Создание элемента
\RequirePackage{enumitem}
\renewcommand{\abbreviation}[2]
{
	\item[---] \noindent #1 -- #2
}

% Создание страницы определений
\renewenvironment{definitions}
{
	\chapter*{ОПРЕДЕЛЕНИЯ}
	\addcontentsline{toc}{chapter}{ОПРЕДЕЛЕНИЯ}
	
	В настоящей расчетно-пояснительной записке применяют следующие термины с соответствующими определениями.
	
	\begin{itemize}[leftmargin=19pt]
	}
	{
	\end{itemize}
}

% Создание страницы обозначений и сокращений
\renewenvironment{abbreviations}
{
	\chapter*{ОБОЗНАЧЕНИЯ И СОКРАЩЕНИЯ}
	\addcontentsline{toc}{chapter}{ОБОЗНАЧЕНИЯ И СОКРАЩЕНИЯ}
	
	В настоящей расчетно-пояснительной записке применяют следующие сокращения и обозначения.
	
	\begin{itemize}[leftmargin=19pt]
	}
	{
	\end{itemize}
}

\usepackage{titlesec} % уже подключается классом, но повторная загрузка не страшна
\titleformat{name=\chapter,numberless}[block]{}{}{0pt}{\large\bfseries\centering}
\titlespacing{name=\chapter,numberless}{0pt}{-22pt}{10pt}